\documentclass[aspectratio=169,11pt]{beamer}

% ============================================
% PACKAGES
% ============================================
\usepackage[utf8]{inputenc}
\usepackage[T1]{fontenc}
\usepackage[french]{babel}
\usepackage{graphicx}
\usepackage{tikz}
\usetikzlibrary{shapes.geometric,shapes.symbols}
\usepackage{pgfplots}
\pgfplotsset{compat=1.18}
\usepackage{booktabs}
\usepackage{listings}
\usepackage{xcolor}
\usepackage{hyperref}
\usepackage{amsmath}

% ============================================
% THEME ET COULEURS
% ============================================
\usetheme{Madrid}
\usecolortheme{whale}

\definecolor{fogblue}{RGB}{37, 99, 235}
\definecolor{foggreen}{RGB}{22, 163, 74}
\definecolor{fogred}{RGB}{220, 38, 38}
\definecolor{fogyellow}{RGB}{202, 138, 4}
\definecolor{codebg}{RGB}{241, 245, 249}

\setbeamercolor{structure}{fg=fogblue}
\setbeamercolor{title}{fg=white,bg=fogblue}
\setbeamercolor{frametitle}{fg=white,bg=fogblue}

% ============================================
% LISTINGS CONFIG
% ============================================
\lstset{
    basicstyle=\ttfamily\scriptsize,
    backgroundcolor=\color{codebg},
    frame=single,
    breaklines=true,
    keywordstyle=\color{fogblue}\bfseries,
    commentstyle=\color{foggreen},
    stringstyle=\color{fogred},
    numbers=none,
    showstringspaces=false,
    extendedchars=true,
    inputencoding=utf8
}

% ============================================
% TITRE
% ============================================
\title[Fog-Assisted Vehicular Computing]{Fog-Assisted Vehicular Fog Computing}
\subtitle{Architecture, Impl\'ementation et D\'emonstration}
\author{Master 2 - Fog \& Edge Computing}
\institute{Universit\'e}
\date{Janvier 2026}

\begin{document}

% ============================================
% PAGE DE TITRE
% ============================================
\begin{frame}
\titlepage
\end{frame}

% ============================================
% PLAN
% ============================================
\begin{frame}{Plan de la Pr\'esentation}
\tableofcontents
\end{frame}

% ############################################
% SECTION 1 : CONTEXTE ET MOTIVATION
% ############################################
\section{Contexte, Motivation et Concepts de Base}

\begin{frame}{Probl\'ematique du Cloud Computing V\'ehiculaire}
\begin{columns}
\column{0.5\textwidth}
\textbf{Limitations du Cloud :}
\begin{itemize}
    \item \textcolor{fogred}{Latence \'elev\'ee} : 100-200 ms
    \item Congestion r\'eseau backbone
    \item D\'ependance \`a la connectivit\'e
    \item Co\^uts de bande passante
    \item Point unique de d\'efaillance
\end{itemize}

\column{0.5\textwidth}
\textbf{Exigences V2X :}
\begin{itemize}
    \item Latence $<$ 20 ms (s\'ecurit\'e)
    \item Latence $<$ 100 ms (infotainment)
    \item Fiabilit\'e 99.999\%
    \item Traitement temps r\'eel
\end{itemize}
\end{columns}

\vspace{0.5cm}
\begin{alertblock}{Constat}
Le Cloud seul ne peut pas r\'epondre aux exigences des applications v\'ehiculaires critiques.
\end{alertblock}
\end{frame}

\begin{frame}{Introduction au Fog Computing}
\begin{definition}
Le \textbf{Fog Computing} est une extension du Cloud qui rapproche le calcul, le stockage et le r\'eseau des utilisateurs finaux, \`a la p\'eriph\'erie du r\'eseau.
\end{definition}

\vspace{0.3cm}
\textbf{Caract\'eristiques cl\'es :}
\begin{itemize}
    \item \textbf{Proximit\'e} : Traitement proche des donn\'ees
    \item \textbf{Faible latence} : 5-50 ms
    \item \textbf{Distribution g\'eographique} : N\oe uds dispers\'es
    \item \textbf{Conscience du contexte} : Localisation, mobilit\'e
    \item \textbf{H\'et\'erog\'en\'eit\'e} : Ressources vari\'ees
\end{itemize}

\vspace{0.3cm}
\begin{block}{Origine}
Terme introduit par Cisco en 2012 (Bonomi et al., "Fog Computing and Its Role in the Internet of Things")
\end{block}
\end{frame}

\begin{frame}{R\'eseaux V\'ehiculaires (VANET) - D\'efinition}
\begin{definition}
\textbf{VANET} (Vehicular Ad-hoc Network) = R\'eseau ad-hoc form\'e par des v\'ehicules communicants, auto-organis\'e et d\'ecentralis\'e.
\end{definition}

\textbf{Caract\'eristiques sp\'ecifiques :}
\begin{itemize}
    \item \textbf{Haute mobilit\'e} : V\'ehicules \`a 50-150 km/h
    \item \textbf{Topologie dynamique} : Liens instables, changements fr\'equents
    \item \textbf{Ressources \'energ\'etiques} : Batterie v\'ehicule (pas de limite)
    \item \textbf{Mod\`eles de mouvement} : Contraints par routes
    \item \textbf{Densit\'e variable} : Zone urbaine vs autoroute
\end{itemize}

\vspace{0.3cm}
\textbf{Standards :}
\begin{itemize}
    \item IEEE 802.11p (WAVE) - Wireless Access in Vehicular Environments
    \item IEEE 1609 (DSRC) - Dedicated Short Range Communications
    \item ETSI ITS-G5 - Standard europ\'een
\end{itemize}
\end{frame}

\begin{frame}{Types de Communication V2X}
\begin{center}
\begin{tikzpicture}[scale=0.85,
    vehicle/.style={draw, circle, fill=fogblue!30, minimum size=1.8cm, font=\small\bfseries, line width=1pt},
    rsu/.style={draw, rectangle, fill=fogyellow!40, minimum width=2cm, minimum height=1cm, font=\small\bfseries, rounded corners=3pt, line width=1pt},
    cloud/.style={draw, ellipse, fill=fogred!25, minimum width=2.2cm, minimum height=1.2cm, font=\small\bfseries, line width=1pt},
    person/.style={draw, circle, fill=foggreen!30, minimum size=1.4cm, font=\scriptsize\bfseries, line width=1pt},
    arrow/.style={->, >=stealth, line width=1.5pt}
]
    % Route (background)
    \fill[gray!20] (-5,-0.6) rectangle (3,0.6);
    \draw[gray!50, dashed] (-5,0) -- (3,0);
    
    % Vehicules sur la route
    \node[vehicle] (v1) at (-3,0) {V\'eh 1};
    \node[vehicle] (v2) at (1,0) {V\'eh 2};
    
    % V2V - entre vehicules
    \draw[arrow, fogblue] (v1) -- node[above, font=\footnotesize\bfseries, fill=white, inner sep=2pt] {V2V} (v2);
    \draw[arrow, fogblue] (v2) -- (v1);
    
    % RSU en hauteur
    \node[rsu] (rsu) at (5,2) {RSU};
    
    % V2I - vehicule vers RSU
    \draw[arrow, fogyellow] (v2.north east) -- node[above, sloped, font=\footnotesize\bfseries, fill=white, inner sep=2pt] {V2I} (rsu.south west);
    \draw[arrow, fogyellow] (rsu.south west) -- (v2.north east);
    
    % Cloud en haut a droite
    \node[cloud] (cloud) at (9,2) {Cloud};
    
    % V2N - RSU vers Cloud
    \draw[arrow, fogred] (rsu) -- node[above, font=\footnotesize\bfseries, fill=white, inner sep=2pt] {V2N} (cloud);
    \draw[arrow, fogred] (cloud) -- (rsu);
    
    % Pieton en bas
    \node[person] (ped) at (-3,-2.5) {Pi\'eton};
    
    % V2P - vehicule vers pieton
    \draw[arrow, foggreen] (v1.south) -- node[left, font=\footnotesize\bfseries, fill=white, inner sep=2pt] {V2P} (ped.north);
    \draw[arrow, foggreen] (ped.north) -- (v1.south);
    
\end{tikzpicture}
\end{center}

\begin{table}
\centering
\scriptsize
\begin{tabular}{lccl}
\toprule
\textbf{Type} & \textbf{Latence} & \textbf{Port\'ee} & \textbf{Application} \\
\midrule
V2V & 1-10 ms & 100-300 m & \'Evitement collision \\
V2I & 10-50 ms & 300-1000 m & Gestion trafic \\
V2N & 50-200 ms & Illimit\'ee & Services Cloud \\
V2P & 1-10 ms & 50 m & S\'ecurit\'e pi\'etons \\
\bottomrule
\end{tabular}
\end{table}
\end{frame}

\begin{frame}{Motivation : Pourquoi le Fog V\'ehiculaire ?}
\begin{center}
\begin{tikzpicture}[scale=0.9]
    \draw[fill=fogred!20, rounded corners] (0,4) rectangle (3,5);
    \node at (1.5,4.5) {\textbf{Cloud}};
    \node at (4.5,4.5) {125 ms - Global};
    
    \draw[fill=fogyellow!20, rounded corners] (0,2) rectangle (3,3);
    \node at (1.5,2.5) {\textbf{Fog (RSU)}};
    \node at (4.5,2.5) {22 ms - Zone};
    
    \draw[fill=foggreen!20, rounded corners] (0,0) rectangle (3,1);
    \node at (1.5,0.5) {\textbf{Edge (V\'ehicule)}};
    \node at (4.5,0.5) {5 ms - Local};
    
    \draw[->, thick] (1.5,3) -- (1.5,4);
    \draw[->, thick] (1.5,1) -- (1.5,2);
    
    \node at (7,4.5) {\textcolor{fogred}{5\% des donn\'ees}};
    \node at (7,2.5) {\textcolor{fogyellow}{25\% des donn\'ees}};
    \node at (7,0.5) {\textcolor{foggreen}{70\% des donn\'ees}};
\end{tikzpicture}
\end{center}

\textbf{B\'en\'efices :} R\'eduction latence 85\%, \'economie bande passante 95\%
\end{frame}

% ############################################
% SECTION 2 : ARCHITECTURE
% ############################################
\section{Architecture Fog-Assisted Vehicular Computing}

\begin{frame}{Architecture 3-Tiers}
\begin{center}
\begin{tikzpicture}[scale=0.8, every node/.style={scale=0.8}]
    \draw[fill=fogblue!10, rounded corners] (-4,5.5) rectangle (4,7);
    \node[font=\bfseries] at (0,6.5) {CLOUD LAYER};
    \node at (0,6) {\small Datacenter, ML global, historique};
    
    \draw[fill=fogyellow!10, rounded corners] (-4,2.5) rectangle (4,5);
    \node[font=\bfseries] at (0,4.5) {FOG LAYER (RSU)};
    \node at (-2,3.5) {RSU\_J1};
    \node at (0,3.5) {RSU\_J2};
    \node at (2,3.5) {RSU\_J3};
    \node at (0,3) {\small Agr\'egation, contr\^ole trafic};
    
    \draw[fill=foggreen!10, rounded corners] (-4,-0.5) rectangle (4,2);
    \node[font=\bfseries] at (0,1.5) {EDGE LAYER (V\'ehicules)};
    \node at (0,0.5) {\small Fog Nodes mobiles, capteurs, traitement local};
    
    \draw[<->, thick, fogblue] (0,5) -- (0,5.5);
    \draw[<->, thick, fogyellow] (-2,2) -- (-2,2.5);
    \draw[<->, thick, fogyellow] (0,2) -- (0,2.5);
    \draw[<->, thick, fogyellow] (2,2) -- (2,2.5);
\end{tikzpicture}
\end{center}
\end{frame}

\begin{frame}[fragile]{Couche Cloud - Configuration}
\textbf{R\^ole :} Traitement global, stockage long terme, Machine Learning

\begin{columns}
\column{0.5\textwidth}
\textbf{Caract\'eristiques :}
\begin{itemize}
    \item Ressources illimit\'ees
    \item Latence : 100-200 ms
    \item Traitement diff\'er\'e acceptable
    \item Vue globale du r\'eseau
\end{itemize}

\column{0.5\textwidth}
\textbf{T\^aches :}
\begin{itemize}
    \item Analyse historique
    \item Entra\^inement mod\`eles ML
    \item Optimisation globale
    \item Backup et archivage
\end{itemize}
\end{columns}

\vspace{0.3cm}
\textbf{Configuration iFogSim :}
\begin{lstlisting}[language=Java]
FogDevice cloud = createFogDevice("cloud", 
    44800,  // MIPS (puissance calcul)
    40000,  // RAM (MB)
    100,    // Bande passante montante (Mbps)
    10000   // Bande passante descendante
);
\end{lstlisting}
\end{frame}

\begin{frame}[fragile]{Couche Fog (RSU) - Sp\'ecifications Techniques}
\textbf{Roadside Units} - Points d'acc\`es fixes aux intersections

\begin{columns}
\column{0.5\textwidth}
\textbf{Sp\'ecifications IEEE 802.11p :}
\begin{itemize}
    \item \textbf{Fr\'equence} : 5.9 GHz (bande ITS)
    \item \textbf{Canaux} : 7 x 10 MHz
    \item \textbf{D\'ebit} : 6-27 Mbps (OFDM)
    \item \textbf{Port\'ee} : 300-1000 m
    \item \textbf{Modulation} : BPSK \`a 64-QAM
\end{itemize}

\column{0.5\textwidth}
\textbf{Dans notre simulation :}
\begin{itemize}
    \item Port\'ee : 300 m
    \item Capacit\'e : 30 v\'ehicules
    \item 4 RSU aux intersections
    \item Latence : 22 ms (base)
\end{itemize}
\end{columns}

\vspace{0.3cm}
\begin{lstlisting}[language=Java]
FogDevice rsu = createFogDevice("RSU_J1",
    2800,   // MIPS
    4000,   // RAM (MB)
    10000,  // Uplink BW
    10000   // Downlink BW
);
\end{lstlisting}
\end{frame}

\begin{frame}[fragile]{Couche Edge (V\'ehicules) - Types}
\textbf{Vehicular Fog Nodes} - Ressources de calcul embarqu\'ees

\begin{table}
\centering
\begin{tabular}{lccc}
\toprule
\textbf{Type} & \textbf{Vitesse max} & \textbf{Fog-capable} & \textbf{MIPS} \\
\midrule
car & 50 km/h & Conditionnel & 500 \\
fog\_capable & 50 km/h & Toujours & 1000 \\
bus & 40 km/h & Toujours & 1500 \\
truck & 35 km/h & Conditionnel & 800 \\
\bottomrule
\end{tabular}
\end{table}

\textbf{Crit\`ere de s\'election Fog Node :}
\begin{lstlisting}[language=Python]
is_fog_node = (
    type == "fog_capable" or  # Equipement dedie
    type == "bus" or          # Arrets frequents = stable
    speed < 8.33              # < 30 km/h = connexion stable
)
\end{lstlisting}

\textbf{Justification :} Un v\'ehicule lent offre une connexion plus stable avec le RSU.
\end{frame}

% ############################################
% SECTION 3 : VEHICULAR FOG NODES
% ############################################
\section{Vehicular Fog Nodes et Fonctionnement Interne}

\begin{frame}{Capacit\'es Embarqu\'ees}
\textbf{Ressources disponibles sur un v\'ehicule moderne :}

\begin{columns}
\column{0.5\textwidth}
\textbf{Mat\'eriel :}
\begin{itemize}
    \item CPU multi-c\oe ur (ARM/x86)
    \item GPU embarqu\'e (NVIDIA Jetson)
    \item RAM : 4-16 GB
    \item Stockage : SSD 128+ GB
    \item Capteurs : LiDAR, cam\'eras, radar
\end{itemize}

\column{0.5\textwidth}
\textbf{Communication :}
\begin{itemize}
    \item DSRC/C-V2X (d\'edi\'e)
    \item 4G/5G cellulaire
    \item WiFi 802.11ac/ax
    \item Bluetooth 5.0
    \item CAN bus interne
\end{itemize}
\end{columns}

\vspace{0.5cm}
\begin{alertblock}{Observation}
Un v\'ehicule moderne dispose de plus de puissance de calcul qu'un datacenter des ann\'ees 2000.
\end{alertblock}
\end{frame}

\begin{frame}{Coop\'eration entre Fog Nodes - Algorithmes}
\textbf{Algorithmes de coop\'eration utilis\'es :}

\begin{enumerate}
    \item \textbf{Task Partitioning} : Division de t\^aches en sous-t\^aches parall\'elisables
    \item \textbf{Load Balancing} : \'Equilibrage dynamique entre v\'ehicules voisins
    \item \textbf{Proximity-Based Clustering} : Formation de clusters locaux
    \item \textbf{Cooperative Caching} : Mise en cache distribu\'ee (LRU)
\end{enumerate}

\vspace{0.3cm}
\textbf{S\'election du meilleur voisin :}
\begin{equation}
\text{score}(v) = \alpha \cdot \frac{1}{\text{charge}_{v}} + \beta \cdot \frac{1}{\text{distance}_{v}} + \gamma \cdot \text{stabilit\'e}_{v}
\end{equation}

O\`u :
\begin{itemize}
    \item $\alpha = 0.4$ (importance de la charge CPU)
    \item $\beta = 0.3$ (importance de la proximit\'e)
    \item $\gamma = 0.3$ (importance de la stabilit\'e de connexion)
\end{itemize}
\end{frame}

\begin{frame}{Coop\'eration - Exemple Concret}
\begin{center}
\begin{tikzpicture}[scale=0.8]
    \node[draw, circle, fill=fogblue!20, minimum size=1.2cm] (A) at (-2,0) {Veh A};
    \node[draw, circle, fill=foggreen!20, minimum size=1.2cm] (B) at (2,0) {Veh B};
    \node[draw, circle, fill=fogyellow!20, minimum size=1.2cm] (C) at (0,2) {Veh C};
    
    \node at (-2,-0.9) {\scriptsize 60 km/h};
    \node at (-2,-1.3) {\scriptsize CPU: 80\%};
    
    \node at (2,-0.9) {\scriptsize 5 km/h};
    \node at (2,-1.3) {\scriptsize CPU: 20\%};
    
    \node at (0,2.9) {\scriptsize 40 km/h};
    \node at (0,3.3) {\scriptsize CPU: 50\%};
    
    \draw[->, thick, fogred] (A) -- node[above, sloped] {\scriptsize T\^ache lourde} (B);
    \draw[->, thick, foggreen, dashed] (B) to[bend left] node[below, sloped] {\scriptsize R\'esultat} (A);
    
    \draw[<->, thick, gray, dashed] (A) -- (C);
    \draw[<->, thick, gray, dashed] (B) -- (C);
\end{tikzpicture}
\end{center}

\textbf{Sc\'enario :} V\'ehicule A (rapide, surcharg\'e) d\'el\`egue vers V\'ehicule B (lent, libre)

\textbf{Gain :} Latence $5$ ms au lieu de $125$ ms (Cloud)
\end{frame}

\begin{frame}{R\^ole des RSU - Fonctions D\'etaill\'ees}
\textbf{Fonctions principales :}

\begin{itemize}
    \item \textbf{Agr\'egation} : Collecte/fusion des donn\'ees de zone (moyenne, max, min)
    \item \textbf{Orchestration} : Affectation des t\^aches aux fog nodes optimaux
    \item \textbf{Cache} : Stockage LRU des donn\'ees fr\'equemment demand\'ees
    \item \textbf{Contr\^ole} : Gestion adaptative des feux de circulation
    \item \textbf{Relais} : Pont vers le Cloud (donn\'ees agr\'eg\'ees seulement)
\end{itemize}

\vspace{0.3cm}
\textbf{Calcul de congestion RSU :}
\begin{equation}
\text{congestion} = \min\left(\frac{N_{\text{v\'ehicules}}}{C_{\text{max}}}, 1.0\right) = \min\left(\frac{N}{30}, 1.0\right)
\end{equation}

\begin{exampleblock}{Seuil de 30 v\'ehicules}
Bas\'e sur la capacit\'e du canal IEEE 802.11p avec protocole CSMA/CA et slot time de 13 $\mu$s.
\end{exampleblock}
\end{frame}

% ############################################
% SECTION 4 : GESTION RESSOURCES
% ############################################
\section{Gestion des Ressources, Mobilit\'e et S\'ecurit\'e}

\begin{frame}{Offloading : Principe et D\'ecision}
\textbf{Offloading} = D\'el\'egation d'une t\^ache vers une ressource externe

\textbf{Crit\`eres de d\'ecision :}
\begin{itemize}
    \item Charge CPU locale ($C_{local}$)
    \item Latence r\'eseau disponible ($L_{net}$)
    \item \'Energie restante batterie ($E$)
    \item Temps d'ex\'ecution estim\'e ($T_{exec}$)
    \item Priorit\'e de la t\^ache ($P$)
\end{itemize}

\vspace{0.3cm}
\textbf{Formule de d\'ecision :}
\begin{equation}
\text{destination} = \arg\min_{d \in \{local, fog, cloud\}} \left( T_d + \alpha \cdot L_d + \beta \cdot E_d \right)
\end{equation}

O\`u $T_d$ = temps d'ex\'ecution, $L_d$ = latence, $E_d$ = co\^ut \'energ\'etique
\end{frame}

\begin{frame}{Calcul d'Offloading - Mod\`ele Math\'ematique}
\textbf{Temps total d'offloading :}
\begin{equation}
T_{offload} = T_{upload} + T_{exec} + T_{download}
\end{equation}

\textbf{D\'etail de chaque composante :}
\begin{align}
T_{upload} &= \frac{D_{input}}{BW_{up}} \\
T_{exec} &= \frac{I}{MIPS_{target}} \\
T_{download} &= \frac{D_{output}}{BW_{down}}
\end{align}

\textbf{Variables :}
\begin{itemize}
    \item $D_{input/output}$ : Taille des donn\'ees (bits)
    \item $BW$ : Bande passante (bps)
    \item $I$ : Instructions requises
    \item $MIPS$ : Millions d'instructions par seconde
\end{itemize}
\end{frame}

\begin{frame}{Calcul des Latences Dynamiques}
\textbf{Mod\`ele utilis\'e dans notre simulation :}

\begin{align}
L_{edge} &= 5.2 + \text{congestion} \times 3 \text{ ms} \\
L_{fog} &= 22.4 + \text{congestion} \times 10 \text{ ms} \\
L_{cloud} &= 125.3 + \text{congestion} \times 25 \text{ ms}
\end{align}

\textbf{Justification des facteurs d'impact :}
\begin{table}
\centering
\begin{tabular}{lcl}
\toprule
\textbf{Niveau} & \textbf{Facteur} & \textbf{Raison} \\
\midrule
Edge ($\times$3) & Faible & Pas de r\'eseau, contention CPU \\
Fog ($\times$10) & Moyen & CSMA/CA, collisions radio \\
Cloud ($\times$25) & \'Elev\'e & Backhaul, Internet, datacenter \\
\bottomrule
\end{tabular}
\end{table}
\end{frame}

\begin{frame}[fragile]{Gestion de la Mobilit\'e - D\'efis}
\textbf{D\'efis li\'es \`a la mobilit\'e v\'ehiculaire :}

\begin{itemize}
    \item \textbf{Handover fr\'equent} : Changement de RSU toutes les 30-60s
    \item \textbf{Topologie dynamique} : V\'ehicules entrent/sortent de zone
    \item \textbf{Continuit\'e de service} : Migration des t\^aches en cours
    \item \textbf{Pr\'ediction} : Anticipation des mouvements futurs
\end{itemize}

\vspace{0.3cm}
\textbf{Solution : Affectation RSU dynamique}
\begin{lstlisting}[language=Python]
def find_nearest_rsu(x, y):
    min_dist = float('inf')
    nearest = None
    for rsu_id, (rx, ry) in RSU_POSITIONS.items():
        dist = sqrt((x-rx)**2 + (y-ry)**2)
        if dist < min_dist and dist < RSU_RANGE:
            min_dist = dist
            nearest = rsu_id
    return nearest
\end{lstlisting}
\end{frame}

\begin{frame}{Gestion de la Mobilit\'e - Strat\'egies}
\textbf{Strat\'egies de handover :}

\begin{enumerate}
    \item \textbf{Proactive Handover} : Pr\'ediction par ML de la trajectoire
    \item \textbf{Pre-caching} : Envoi anticipatif des donn\'ees vers RSU suivant
    \item \textbf{Task Migration} : D\'eplacement des t\^aches actives
    \item \textbf{Session Persistence} : Maintien de l'\'etat entre RSUs
\end{enumerate}

\vspace{0.3cm}
\textbf{M\'etriques de d\'ecision :}
\begin{equation}
\text{handover\_trigger} = \text{RSSI}_{current} < \text{RSSI}_{threshold} + \text{hystereris}
\end{equation}

\begin{itemize}
    \item RSSI threshold = -85 dBm (typique 802.11p)
    \item Hyst\'er\'esis = 3 dB (\'evite ping-pong)
\end{itemize}
\end{frame}

\begin{frame}{S\'ecurit\'e et Protection des Donn\'ees}
\textbf{Menaces potentielles dans les VANET :}

\begin{columns}
\column{0.5\textwidth}
\textbf{Attaques :}
\begin{itemize}
    \item \textbf{Sybil Attack} : Fausses identit\'es multiples
    \item \textbf{Man-in-the-Middle} : Interception des messages
    \item \textbf{DoS/DDoS} : Saturation du canal
    \item \textbf{Replay Attack} : R\'eenvoi de messages anciens
    \item \textbf{False Information} : Faux messages de s\'ecurit\'e
\end{itemize}

\column{0.5\textwidth}
\textbf{Contre-mesures :}
\begin{itemize}
    \item \textbf{PKI/Certificats} : IEEE 1609.2
    \item \textbf{Pseudonymes rotatifs} : Vie priv\'ee
    \item \textbf{V\'erification plausibilit\'e} : Coh\'erence des donn\'ees
    \item \textbf{Timestamp/Nonce} : Anti-replay
    \item \textbf{D\'etection anomalies} : ML
\end{itemize}
\end{columns}
\end{frame}

\begin{frame}{S\'ecurit\'e - Standards et Protocoles}
\textbf{Standards de s\'ecurit\'e V2X :}

\begin{table}
\centering
\scriptsize
\begin{tabular}{lll}
\toprule
\textbf{Standard} & \textbf{Fonction} & \textbf{M\'ecanisme} \\
\midrule
IEEE 1609.2 & S\'ecurit\'e messages & ECDSA, AES-CCM \\
ETSI TS 102 940 & Architecture PKI & Certificats X.509 \\
ETSI TS 102 941 & Gestion certificats & CRL, OCSP \\
SAE J2945 & Formats messages & BSM sign\'es \\
\bottomrule
\end{tabular}
\end{table}

\vspace{0.3cm}
\textbf{Exemple signature ECDSA :}
\begin{itemize}
    \item Courbe : NIST P-256
    \item Taille signature : 64 octets
    \item Temps signature : $\sim$1 ms
    \item Temps v\'erification : $\sim$2 ms
\end{itemize}

\begin{alertblock}{Note}
Notre simulation actuelle ne mod\'elise pas la s\'ecurit\'e - focus sur les performances r\'eseau.
\end{alertblock}
\end{frame}

% ############################################
% SECTION 5 : CAS PRATIQUE
% ############################################
\section{Cas Pratique : Notre Simulation}

\begin{frame}{Crit\`eres de Choix des Outils}
\textbf{Pourquoi SUMO + Python + iFogSim ?}

\begin{table}
\centering
\scriptsize
\begin{tabular}{llp{5cm}}
\toprule
\textbf{Outil} & \textbf{Alternative} & \textbf{Raison du choix} \\
\midrule
\textbf{SUMO} & NS-3, OMNeT++ & Sp\'ecialis\'e trafic, open-source, TraCI \\
\textbf{Python} & Java, C++ & Rapidit\'e dev, TraCI bindings, Flask \\
\textbf{iFogSim} & EdgeCloudSim, YAFS & Mod\`ele 3-tiers natif, CloudSim base \\
\textbf{Flask} & FastAPI, Django & L\'eger, WebSocket facile \\
\textbf{Chart.js} & D3.js, Plotly & Simple, temps r\'eel ok \\
\bottomrule
\end{tabular}
\end{table}

\textbf{Crit\`eres principaux :}
\begin{itemize}
    \item \textbf{Int\'egrabilit\'e} : SUMO $\leftrightarrow$ Python $\leftrightarrow$ Java
    \item \textbf{R\'ealisme} : SUMO certifi\'e pour simulation trafic
    \item \textbf{Extensibilit\'e} : iFogSim modifiable (open-source)
    \item \textbf{Visualisation} : Dashboard temps r\'eel
\end{itemize}
\end{frame}

\begin{frame}[fragile]{Configuration SUMO - Fichiers}
\textbf{Fichiers de configuration cr\'e\'es :}

\begin{lstlisting}[language=XML]
<!-- simulation.sumocfg -->
<configuration>
    <input>
        <net-file value="network.net.xml"/>
        <route-files value="routes.rou.xml"/>
        <additional-files value="detectors.add.xml"/>
    </input>
    <time>
        <begin value="0"/>
        <end value="900"/>
        <step-length value="1"/>
    </time>
</configuration>
\end{lstlisting}

\textbf{3 phases de simulation :}
\begin{itemize}
    \item 0-300s : Trafic normal ($\sim$50 v\'ehicules)
    \item 300-600s : Congestion ($\sim$150 v\'ehicules)
    \item 600-900s : R\'ecup\'eration ($\sim$80 v\'ehicules)
\end{itemize}
\end{frame}

\begin{frame}{Architecture Logicielle}
\begin{center}
\begin{tikzpicture}[scale=0.75, every node/.style={scale=0.75}]
    \node[draw, rounded corners, fill=fogblue!20, minimum width=3cm, minimum height=1cm] (sumo) at (-4,0) {SUMO-GUI};
    
    \node[draw, rounded corners, fill=foggreen!20, minimum width=3cm, minimum height=1cm] (python) at (0,0) {server.py};
    
    \node[draw, rounded corners, fill=fogyellow!20, minimum width=3cm, minimum height=1cm] (ifogsim) at (4,0) {iFogSim Java};
    
    \node[draw, rounded corners, fill=fogred!20, minimum width=3cm, minimum height=1cm] (dashboard) at (0,-2.5) {Dashboard Web};
    
    \draw[<->, thick] (sumo) -- node[above] {\small TraCI} (python);
    \draw[<->, thick] (python) -- node[above] {\small Socket TCP} (ifogsim);
    \draw[<->, thick] (python) -- node[right] {\small WebSocket} (dashboard);
    
    \node at (-4,-0.7) {\footnotesize Port: 8813};
    \node at (0,-0.7) {\footnotesize Port: 5000};
    \node at (4,-0.7) {\footnotesize Port: 5555};
\end{tikzpicture}
\end{center}

\textbf{Flux de donn\'ees :}
\begin{enumerate}
    \item SUMO g\'en\`ere positions v\'ehicules (TraCI)
    \item Python enrichit et envoie \`a iFogSim (JSON/Socket)
    \item iFogSim calcule m\'etriques fog et renvoie
    \item Dashboard affiche en temps r\'eel (WebSocket)
\end{enumerate}
\end{frame}

\begin{frame}[fragile]{D\'emonstration}
\textbf{Lancement :}

\begin{lstlisting}[language=bash]
# Terminal 1 - iFogSim
cd simulation/ifogsim
java -cp "out;jars/*" org.fog.test.perfeval.VehicularFogSimulation

# Terminal 2 - Python + SUMO
cd simulation
python server.py --dashboard --gui
\end{lstlisting}

\vspace{0.3cm}
\textbf{Acc\`es Dashboard :} \url{http://localhost:5000}

\begin{block}{R\'esultats observ\'es}
\begin{itemize}
    \item Latence moyenne : $\sim$16-20 ms (vs 125 ms Cloud-only)
    \item Fog nodes actifs : 50-80 en phase congestion
    \item R\'eduction bande passante vers Cloud : 95\%
\end{itemize}
\end{block}
\end{frame}

% ############################################
% SECTION 6 : AVANTAGES ET LIMITES
% ############################################
\section{Avantages, Limites et Comparaison}

\begin{frame}{Comparaison Cloud vs Fog-Assisted}
\begin{table}
\centering
\begin{tabular}{lcc}
\toprule
\textbf{M\'etrique} & \textbf{Cloud-Only} & \textbf{Fog-Assisted} \\
\midrule
Latence moyenne & 125 ms & 16-20 ms \\
R\'eduction latence & - & \textcolor{foggreen}{85\%} \\
Bande passante Cloud & 100\% & 5\% \\
\'Economie BP & - & \textcolor{foggreen}{95\%} \\
Scalabilit\'e locale & Faible & \'Elev\'ee \\
Tol\'erance aux pannes & Faible & \'Elev\'ee \\
\bottomrule
\end{tabular}
\end{table}

\vspace{0.3cm}
\textbf{Formule latence moyenne pond\'er\'ee :}
\begin{equation}
L_{avg} = \frac{L_{edge} \times 70 + L_{fog} \times 25 + L_{cloud} \times 5}{100}
\end{equation}
\end{frame}

\begin{frame}{Avantages du Fog Computing V\'ehiculaire}
\begin{columns}
\column{0.5\textwidth}
\textbf{Performance :}
\begin{itemize}
    \item[$\checkmark$] Faible latence ($<$ 20 ms)
    \item[$\checkmark$] Traitement temps r\'eel
    \item[$\checkmark$] R\'eduction congestion r\'eseau
\end{itemize}

\vspace{0.3cm}
\textbf{R\'esilience :}
\begin{itemize}
    \item[$\checkmark$] Pas de point unique de d\'efaillance
    \item[$\checkmark$] Fonctionne hors-ligne (mode d\'egrad\'e)
    \item[$\checkmark$] Distribution g\'eographique
\end{itemize}

\column{0.5\textwidth}
\textbf{\'Economie :}
\begin{itemize}
    \item[$\checkmark$] R\'eduction co\^uts Cloud
    \item[$\checkmark$] Utilisation ressources existantes
    \item[$\checkmark$] Moins de bande passante
\end{itemize}

\vspace{0.3cm}
\textbf{Vie priv\'ee :}
\begin{itemize}
    \item[$\checkmark$] Donn\'ees trait\'ees localement
    \item[$\checkmark$] Moins de donn\'ees vers Cloud
\end{itemize}
\end{columns}
\end{frame}

\begin{frame}{Limites et D\'efis}
\begin{columns}
\column{0.5\textwidth}
\textbf{Techniques :}
\begin{itemize}
    \item[$\times$] Ressources limit\'ees par v\'ehicule
    \item[$\times$] Mobilit\'e = instabilit\'e
    \item[$\times$] H\'et\'erog\'en\'eit\'e des \'equipements
    \item[$\times$] Consommation \'energ\'etique
\end{itemize}

\column{0.5\textwidth}
\textbf{Organisationnels :}
\begin{itemize}
    \item[$\times$] D\'eploiement infrastructure RSU
    \item[$\times$] Standardisation en cours
    \item[$\times$] Mod\`ele \'economique incertain
    \item[$\times$] S\'ecurit\'e complexe
\end{itemize}
\end{columns}

\vspace{0.5cm}
\begin{alertblock}{D\'efi principal}
Comment garantir la QoS avec des ressources volatiles et une topologie dynamique ?
\end{alertblock}
\end{frame}

% ############################################
% SECTION 7 : PERSPECTIVES
% ############################################
\section{Perspectives et Conclusion}

\begin{frame}{Int\'egration 5G/6G}
\textbf{\'Evolution technologique :}

\begin{table}
\centering
\begin{tabular}{lccc}
\toprule
\textbf{Param\`etre} & \textbf{DSRC} & \textbf{5G C-V2X} & \textbf{6G} \\
\midrule
Latence & 20-50 ms & 1-10 ms & $<$ 1 ms \\
D\'ebit & 27 Mbps & 1 Gbps & 1 Tbps \\
Port\'ee & 300 m & 500 m & 1 km \\
Capacit\'e & 30 v\'eh & 100 v\'eh & 1000 v\'eh \\
\bottomrule
\end{tabular}
\end{table}

\vspace{0.3cm}
\end{frame}

\begin{frame}{Smart Cities et V\'ehicules Autonomes}
\textbf{Applications futures :}

\begin{itemize}
    \item \textbf{Conduite autonome niveau 5} : Fog Computing essentiel
    \item \textbf{Gestion trafic intelligente} : Feux adaptatifs via RSU
    \item \textbf{Parking intelligent} : Guidage temps r\'eel
    \item \textbf{Platooning} : Convois de camions autonomes
    \item \textbf{Emergency services} : Priorit\'e automatique
\end{itemize}

\vspace{0.3cm}
\textbf{IA distribu\'ee :}
\begin{itemize}
    \item Federated Learning sur v\'ehicules
    \item Mod\`eles ML l\'egers (TinyML)
    \item Inf\'erence au bord
\end{itemize}
\end{frame}

\begin{frame}{Conclusion}
\textbf{Ce que nous avons d\'emontr\'e :}

\begin{enumerate}
    \item Le \textbf{Fog Computing} r\'epond aux exigences de latence v\'ehiculaire
    \item L'architecture \textbf{3-tiers} (Edge/Fog/Cloud) optimise les ressources
    \item Notre simulation \textbf{SUMO + iFogSim} valide le concept
    \item R\'eduction de latence de \textbf{85\%} par rapport au Cloud seul
\end{enumerate}

\vspace{0.5cm}
\begin{block}{Message cl\'e}
Le Fog Computing V\'ehiculaire n'est pas une alternative au Cloud, mais un \textbf{compl\'ement essentiel} pour les applications temps r\'eel.
\end{block}
\end{frame}

\begin{frame}{R\'ef\'erences}
\footnotesize
\begin{thebibliography}{9}
\bibitem{bonomi2012} Bonomi, F. et al. (2012). "Fog Computing and Its Role in the Internet of Things." MCC Workshop.

\bibitem{hou2016} Hou, X. et al. (2016). "Vehicular Fog Computing: A Viewpoint of Vehicles as the Infrastructures." IEEE VTC.

\bibitem{ifogsim} Gupta, H. et al. (2017). iFogSim: A toolkit for modeling and simulation of resource management techniques in the Internet of Things, Edge and Fog computing environments.

\end{thebibliography}
\end{frame}

\begin{frame}
\centering
\Huge\textbf{Merci pour votre attention}

\vspace{1cm}
\Large Questions ?

\vspace{1cm}
\end{frame}

\end{document}
